\documentclass[11pt,a4paper]{article}
\usepackage[utf8]{inputenc}
\usepackage[T1]{fontenc}
\usepackage[spanish]{babel}
\usepackage{hyperref}
\usepackage{geometry}
\geometry{margin=2.5cm}
\usepackage{parskip}

\title{002 procesos}
\author{Alejandro Valenzuela Rivera\\\small Universidad Autónoma de Nuevo León}
\date{\today}
\begin{document}
\maketitle

\begin{quote}
\textbf{Sesgo declarado.} Este trabajo se sostiene desde el sesgo declarado de su autor, y no finge neutralidad. Lo que afirma se puede tumbar, y ese es el punto: el falsificador está escrito, no escondido.
\end{quote}

\section*{Resumen}
Procesos es el primitivo de lo que se repite: entra, falla, sale y consume.

\section{Por qué se sostiene}
Lo que vuelve se puede correr sin su autor.

\section{De dónde viene}
La autoridad, en las conversaciones de su historial de sesgos.

\section{El caso}
Ninguna: es la forma de lo que se repite.

\section{Método}
Un proceso que corre igual sin su autor, y que aun asi cambia algo.

\section{Límites, y qué lo tumba}
Un procedimiento que solo funciona si el que lo hizo esta ahi.

\section{Lo que queda abierto}
\textquestiondown{}Que hace que una posicion sobreviva a su autor?

\section*{Cómo se leyó esta tesis}
Esto se escribió en el repositorio man-made-empathy, del que este candado es una pieza: candados/01 el ente y el metodo/002 procesos.md. La máquina que lo compuso lo hizo desde este mismo archivo, y no agregó nada que no estuviera escrito. Si algo de esto es falso, se dice en el mismo sitio, y se corrige con fecha.

\end{document}
